\chapter{Protokol zberu údajov na slovenskej škole}
\label{chapter:protokol}

Táto príloha je návod na vykonanie, nie opis zámeru. Je napísaná tak, aby podľa nej mohol zber viesť učiteľ informatiky po jednom zaškolení, a aby sa dalo skontrolovať, či prebehol podľa plánu.

\section{Čo škola poskytuje a čo dostane}

\begin{description}[leftmargin=0.6cm, nosep]
    \item[\textbf{Čo škola dáva.}] Tri vyučovacie bloky na výučbu (90~min každý, v rámci bežných hodín informatiky), \asmMeasureLessons{} vyučovacie hodiny na meranie, počítačovú učebňu s pripojením a jedného učiteľa, ktorý vedie všetky tri bloky v danej triede.
    \item[\textbf{Čo škola nedáva.}] Zásah do siete, inštaláciu serverov, prístup k svojim systémom ani osobné údaje žiakov. Nástroj beží lokálne na notebookoch a po skončení sa odstráni vrátane certifikačnej autority.
    \item[\textbf{Čo škola dostane.}] Tri pripravené hodiny informatiky s materiálmi, ktoré si ponecháva; agregovanú správu za triedu (nie za žiaka) s porovnaním voči ostatným zapojeným triedam; a po skončení štúdie prístup k celej sade úloh a testov.
\end{description}

\section{Harmonogram a jeho miesto v školskom roku}

\begin{table}[H]
\centering
\caption{Harmonogram jednej triedy}
\label{tab:harmonogram}
\small
\begin{tabularx}{\textwidth}{@{}l l l X@{}}
\toprule
\textbf{Bod} & \textbf{Kedy} & \textbf{Trvanie} & \textbf{Obsah} \\
\midrule
S0   & 4 týždne vopred & --- & súhlasy domov cez EduPage; súhlas riaditeľa; potvrdenie termínov \\
T0a  & týždeň 0 & \mplMinToa~min z \mplBudgetToa & ICAR (\mplGfItems{} položiek), iBobor (8 úloh) \\
T0b  & týždeň 0 & \mplMinTob~min z \mplBudgetTob & Jr.~MAI a vstupný dotazník, \emph{cez EduPage ako domáca úloha} \\
P1--P3 & týždne 1--3 & 2 hodiny & hodina (\instrLessonMin~min), okamžitý test s JOL (\asmFormItems), dotazník \\
T4a  & týždeň 6 & 2 hodiny & odložený test (3 $\times$ \asmFormItems), ICAR znovu (\mplGfItems) \\
T4b  & týždeň 6 & 1 hodina & transferový test (\instrTtItems), záverečný dotazník \\
T5   & týždeň 7 & 30~min/žiak & rozhovory s 18~žiakmi vybranými podľa tercilu PEQI \\
T6   & voliteľne, koniec semestra & \mplMinTsix~min z \mplBudgetTsix & iBobor (15 úloh) ako výstup semestrálneho ramena \\
\bottomrule
\end{tabularx}
\end{table}

\paragraph{Prečo takto a nie inak.} Rozvrh vychádza z troch úprav, z ktorých každá je metodickým zlepšením, nie ústupkom:

\begin{enumerate}[nosep]
    \item \textbf{T0 sa delí.} Dotazník o návykoch nie je test a nepotrebuje dozor. Odchádza cez EduPage ako domáca úloha a prestáva súťažiť o hodinu s nástrojmi, ktoré merajú na čas.
    \item \textbf{Vstupné vedomosti sa presúvajú.} Namiesto jedného dvadsaťpoložkového testu v týždni~0 sa zadá päťpoložkový predtest k \emph{danej} téme, a to na konci predchádzajúcej hodiny informatiky. Globálny test tri týždne pred treťou témou je horší kovariát než päť položiek zadaných pár dní pred hodinou, ktorá ich potrebuje, a stojí štvornásobok času. Zadanie mimo vyučovacieho bloku je nutnosť: blok na päť položiek nemá miesto (príloha~\ref{chapter:hodnotiaci-modul}, časť~\ref{sec:cas-rozpocet}).
    \item \textbf{Blok iBobor sa v hlavnom ramene skracuje.} Tam je kovariátom, nie výstupom. V plnej dĺžke sa vracia len vo voliteľnom semestrálnom ramene, kde \emph{je} výstupom a kde horizont merania konečne zodpovedá rýchlosti zmeny konštruktu (časť~\ref{sec:konstrukty}).
\end{enumerate}

\paragraph{Umiestnenie v školskom roku.} Zber sa plánuje na október až november alebo na marec až apríl. Vylučujú sa dva týždne pred polročnou a koncoročnou klasifikáciou, lebo vtedy je ochota škôl uvoľniť hodiny najnižšia a pozornosť žiakov je viazaná na známkovanie. Termín odloženého testu sa žiakom neoznamuje; učiteľ ho má v pláne ako „opakovaciu hodinu“.

\section{Role}

\begin{description}[leftmargin=0.6cm, nosep]
    \item[\textbf{Učiteľ}] vedie všetky tri bloky v svojej triede. Je to tá istá osoba vo všetkých troch podmienkach~--- to je súčasť kontroly, nie organizačná náhoda. Dostáva scenár hodiny a kontrolný list vernosti.
    \item[\textbf{Výskumník}] je prítomný pri T0a, T4a a T4b, spúšťa záchytnú vrstvu, vypĺňa kontrolný list vernosti a nezasahuje do výučby.
    \item[\textbf{Kontaktná osoba školy}] (spravidla zástupca riaditeľa) zbiera súhlasy, drží soľ pre pseudonymizáciu a je jediná osoba, ktorá vie priradiť kód k žiakovi.
\end{description}

\section{Priebeh jednotlivých stretnutí}

\subsection*{T0a --- vstupná hodina (\mplBudgetToa~min)}

\begin{description}[leftmargin=1.5cm, style=nextline, nosep]
    \item[0--5 min] Rozdanie kódov. Každý žiak dostane obálku so školským kódom. Kód nie je meno ani rodné číslo; priradenie drží kontaktná osoba školy.
    \item[5--8 min] Vysvetlenie: čo sa meria, čo sa nezaznamenáva, že výsledok nemá vplyv na známku a že žiak môže kedykoľvek prestať bez následku. Overenie, že všetci prítomní majú súhlas; žiaci bez súhlasu absolvujú hodinu, ale ich dáta sa nezaznamenávajú.
    \item[8--23 min] ICAR, \mplGfItems{} položiek. Na obrazovke, jedna položka naraz, bez návratu späť. Bez časového limitu na položku, celkový limit 15~minút.
    \item[23--36 min] iBobor, 8 úloh. Návrat k predchádzajúcim úlohám povolený, ako v súťaži.
    \item[36--45 min] Rezerva. Nepoužitá rezerva je zámer: technický problém pri prvom stretnutí je pravidlo, nie výnimka.
\end{description}

\subsection*{T0b --- zadané cez EduPage (\mplBudgetTob~min doma)}

Jr.~MAI (\instrMaiItems~položiek) a vstupný dotazník o doterajšom používaní nástrojov umelej inteligencie a o programátorskej skúsenosti. Zadáva sa ako domáca úloha s týždňovým termínom. Nevyplnenie nie je dôvodom na vylúčenie; chýbajúca hodnota kovariátu sa dopĺňa viacnásobnou imputáciou.

\subsection*{P1--P3 --- vyučovací blok (2 hodiny, 90~min)}

Predtest k téme sa nezadáva v~tomto bloku, ale na konci \emph{predchádzajúcej} hodiny informatiky (4~min, pod dozorom, 2--3 dni vopred). Vedomosti o~predpokladoch témy sa za dva dni bez výučby nemenia, takže kovariát tým nič nestráca, a blok získa štyri minúty, ktoré v~ňom nie sú.

\begin{description}[leftmargin=1.5cm, style=nextline, nosep]
    \item[0--4 min] Spustenie prostredia. V podmienkach $G$ a $A$ sa spúšťa cvičné rozhranie s príslušným prepínačom; v podmienke $K$ to isté rozhranie bez asistenta. Výskumník overí, že sa zapisuje záchytný súbor.
    \item[4--68 min] Hodina podľa scenára, \instrLessonMin{}~minút. Učiteľ vedie rovnakú štruktúru vo všetkých troch podmienkach: výklad, riadené cvičenie, samostatné cvičenie.
    \item[68--84 min] Okamžitý test, \asmFormItems{}~položiek s posuvníkom sebadôvery pri každej uzavretej položke.
    \item[84--89 min] Krátky dotazník: záťaž, pocit naučenia, spokojnosť. Uzatvorenie záchytného súboru a jeho prenos.
\end{description}

\subsection*{T4a --- odložené meranie a negatívna kontrola (2 hodiny)}

Poradie je zámerné: odložený test ide prvý, kým je pozornosť najvyššia, a ICAR posledný, lebo je na únavu najodolnejší. Medzi témami sú trojminútové prestávky, ktoré držia najdlhší nepretržitý úsek testovania pod \instrFatigueCap{}~minútami~--- tridsaťšesť položiek za sebou je jedno takmer päťdesiatminútové sedenie a výkon by v~ňom klesal, a klesal by nerovnako podľa motivácie, ktorú podmienky ovplyvňujú.

\begin{description}[leftmargin=1.5cm, style=nextline, nosep]
    \item[0--5 min] Rozdanie kódov, bez upozornenia vopred.
    \item[5--21 min] Odložený test, téma T1, \asmFormItems{}~položiek, s posuvníkom sebadôvery.
    \item[21--24 min] Prestávka.
    \item[24--40 min] Odložený test, téma T2.
    \item[40--43 min] Prestávka.
    \item[43--59 min] Odložený test, téma T3.
    \item[59--62 min] Prestávka.
    \item[62--77 min] ICAR znovu, \mplGfItems{} položiek~--- tých istých, v tom istom poradí.
\end{description}

\subsection*{T4b --- transferový test (1 hodina)}

Oddelený od T4a. Spolu by to bolo 92~minút merania, a transfer je práve to meranie, ktorému únava uškodí najviac: \instrTtItems{}~rubrikových položiek vyžaduje odvodenie, nie vybavenie.

\begin{description}[leftmargin=1.5cm, style=nextline, nosep]
    \item[0--5 min] Rozdanie kódov.
    \item[5--30 min] Transferový test, \instrTtItems{}~položiek, rubrika 0--3.
    \item[30--40 min] Záverečný dotazník postojov; poďakovanie; informácia, kedy a v akej podobe dostane škola správu.
\end{description}

\paragraph{Čo tieto zmeny stoja a prečo stoja za to.} Počet meracích hodín, ktoré škola dáva, stúpa na \asmMeasureLessons. Pôvodný rozvrh sa zmestil len preto, že časový rozpočet počítal plochý priemer sekúnd na položku a tým podhodnotil rubrikové nástroje o~\asmBudgetUnderstated{}~minúty, najviac transferový test o~\asmBudgetWorstBy{}~minúty (príloha~\ref{chapter:hodnotiaci-modul}, časť~\ref{sec:cas-rozpocet}). Hodina viac je lacnejšia než meranie, ktoré sa do hodiny nevojde a skončí sa predčasne.

\section{Údaje: čo vzniká, kde a v akej podobe}

\begin{table}[H]
\centering
\caption{Toky údajov}
\label{tab:toky}
\small
\begin{tabularx}{\textwidth}{@{}l l X@{}}
\toprule
\textbf{Zdroj} & \textbf{Formát} & \textbf{Čo obsahuje a čo nie} \\
\midrule
testy a dotazníky & CSV & pseudonym, kód položky, odpoveď, čas; \emph{žiadne} meno ani trieda po mene \\
cvičné rozhranie & JSON Lines & udalosti hodiny: zobrazenie úlohy, pokus, prompt, odomknutie \\
záchytná vrstva & JSON Lines & prompty odoslané modelu, po odstránení osobných údajov \\
kontrolný list vernosti & PDF/CSV & či hodina prebehla podľa scenára; vypĺňa výskumník \\
\bottomrule
\end{tabularx}
\end{table}

Pseudonym je školský kód hašovaný soľou, ktorú drží škola. Výskumník teda nemá prostriedok na spätné priradenie, a to ani pri úniku celého analytického súboru. Voľný text sa zaznamenáva len tam, kde to súhlas výslovne pokrýva: prompt a vlastné vysvetlenie. Nikde inde.

\subsection*{Kontrola údajov v deň zberu}

Schéma údajov je strojovo čitateľná a kontroluje sa modulom \texttt{ingest.py} \emph{v ten istý deň}, keď zber prebehol. Dôvod je vecný a v školskom výskume rozhodujúci: chýbajúci stĺpec, preklep v kóde alebo nedoručený súhlas sa pri kontrole až pri analýze objavia o niekoľko mesiacov, keď je škola nedostupná a žiaci sú inde. Vtedy sa z opraviteľnej administratívnej chyby stáva trvalo chýbajúci údaj. Telefonát v deň zberu ju odstráni; marcová analýza už nie.

Kontrola rozlišuje dve triedy nálezov, a to rozlíšenie je zámerné:

\begin{description}[leftmargin=0.6cm, nosep]
    \item[\textbf{Chyby}] bránia analýze a vyžadujú zásah. Patria sem: odpoveď zaznamenaná bez úplného súhlasu (taký riadok sa maže, neanalyzuje); identifikátor, ktorý nie je 16~hexadecimálnych znakov, teda nebol hašovaný a súbor nesmie opustiť školu; pseudonym, ktorý nie je v zozname; názov nástroja mimo schémy; tá istá položka zodpovedaná dvakrát; posuvník sebadôvery na položke dotazníka, ktorá nemá správnu odpoveď; viac vykonaných krokov hodiny než plánovaných.
    \item[\textbf{Upozornenia}] analýzu neblokujú, ale menia zaobchádzanie s blokom: vernosť pod 90\,\%, pokrytie záchytnej vrstvy pod 80\,\%, návratnosť súhlasov pod 85\,\%. Súbor je v poriadku; problém je v tom, čo opisuje.
\end{description}

Prahy sú definované na jednom mieste v kóde, takže protokol a kontrola sa nemôžu rozísť. Kontrola je sama testovaná: k schéme existuje jedenásť zámerne pokazených vzoriek a test vyžaduje, aby bola každá z nich odhalená a aby bezchybná vzorka prešla bez jedinej poznámky. Kontrola, ktorá videla len správne dáta, nie je kontrola, ale tvrdenie.

\section{Súhlas a právny základ}

\begin{description}[leftmargin=0.6cm, nosep]
    \item[\textbf{Žiaci do 18~rokov}] potrebujú písomný informovaný súhlas zákonného zástupcu. Súhlas sa distribuuje cez EduPage a v papierovej forme pre rodiny bez prístupu; obe cesty sa ponúkajú súčasne.
    \item[\textbf{Žiaci nad 18~rokov}] dávajú súhlas sami.
    \item[\textbf{Žiak sám}] dáva súhlas vždy, bez ohľadu na vek, a jeho nesúhlas je rozhodujúci aj vtedy, keď zákonný zástupca súhlasil. Odmietnutie nemá žiadny dôsledok pre známku ani pre účasť na hodine.
    \item[\textbf{Právny základ}] spracúvania je súhlas podľa zákona č.~18/2018~Z.\,z. a nariadenia GDPR. Súhlas je odvolateľný kedykoľvek a bez udania dôvodu; po odvolaní sa záznamy daného pseudonymu vymažú do 30~dní.
    \item[\textbf{Schválenia}] pred začiatkom: etická komisia fakulty, riaditeľ každej zapojenej školy, a súhlas organizátorov súťaže iBobor s použitím úloh na výskumný účel.
\end{description}

\section{Čo robiť, keď niečo zlyhá}

Zoznam je krátky, lebo je vopred rozhodnutý. Improvizácia počas zberu je to, čo robí dáta neporovnateľnými.

\begin{description}[leftmargin=0.6cm, nosep]
    \item[\textbf{Vypadne internet.}] Podmienka $K$ pokračuje bez zmeny. Podmienky $G$ a $A$ sa prerušia, hodina sa vedie do konca bez asistenta a \emph{celý blok danej triedy sa vyradí z analýzy podľa protokolu}; vykazuje sa medzi odchýlkami. Nenahrádza sa iným termínom, lebo perióda je súčasťou dizajnu.
    \item[\textbf{Žiak chýba na jednej perióde.}] Zostáva v štúdii. Krížový dizajn je na neúplné bloky odolný; chýbajúca perióda sa ošetruje zmiešaným modelom.
    \item[\textbf{Žiak chýba na T4.}] Toto je jediný výpadok, ktorý skutočne bolí (časť o citlivosti na výpadok). Náhradný termín sa ponúka do jedného týždňa a odstup sa zaznamená ako kovariát.
    \item[\textbf{Chýba súhlas.}] Žiak absolvuje hodinu ako ktorýkoľvek iný; nezaznamenáva sa oňho nič. Počet takých žiakov sa vykazuje.
    \item[\textbf{Učiteľ sa odchýli od scenára.}] Zaznamená sa v kontrolnom liste. Bloky s vernosťou pod 90\,\% vstupujú do citlivostnej analýzy osobitne.
    \item[\textbf{Žiak použije vlastný telefón.}] Zachytí sa ako nesúlad medzi počtom promptov v rozhraní a v záchytnej vrstve. Vykazuje sa ako miera pokrytia; prípady sa označia a beží citlivostná analýza bez nich.
\end{description}

\section{Kontrolný list pred začiatkom zberu}

\begin{enumerate}[nosep]
    \item Etická komisia schválila, súhlas riaditeľa podpísaný, súhlas organizátorov iBoboru získaný.
    \item Súhlasy vrátené aspoň od 85\,\% žiakov v triede; pod touto hranicou sa trieda nezaraďuje.
    \item Pilot na jednej nezaradenej triede prebehol; položky vyradené podľa \emph{intervalového} pravidla (vyradiť len ak 90\,\% interval leží celý mimo pásma $0{,}30$--$0{,}80$) a nahradené z rezervnej sady. Bodové pravidlo sa nepoužíva: pri \pilotN{}~žiakoch by vyradilo $\pilotPointGood$\,\% dobrých položiek (časť~\ref{sec:pilot}).
    \item Soľ pre pseudonymizáciu vygenerovaná a odovzdaná kontaktnej osobe školy; výskumník ju nemá.
    \item Cvičné rozhranie otestované na školskom notebooku vrátane výpadku siete.
    \item Vrstva podmienky $A$ prešla sadou 40~útočných dopytov s nulovým únikom.
    \item Scenáre hodín a kontrolné listy vernosti vytlačené; učiteľ zaškolený.
    \item Záložný papierový variant testov pripravený pre prípad výpadku.
\end{enumerate}
